\chapter{Motivating Numerical Methods for BVPs}

Up to this point in the course we primarily focused on the 
solution of IVPs, i.e., ODEs with initial conditions 
prescribed at a given point in time which we then solved 
using time-marching methods. In this sense, IVPs have a 
natural interpretation as dynamic phenomenon being simulated
forward in time. We now turn our attention to BVPs, i.e., 
problem of the form 
\begin{equation*}
    \frac{d^n}{dt^n} u = f  \left(t, u, \frac{d u }{dt}, 
    \dots, \frac{d^{n-1} u }{dt^{n-1}} \right), \quad a < t < b, 
\end{equation*}
subject to boundary conditions 
\begin{equation*}
   {\mbf g}(u, a, b) = 0.
\end{equation*}
We will mostly focus on particular boundary conditions shortly. 
The difficulty with BVPs is that their analytic solutions 
are in general difficult to write down and they differ from 
IVPs in nature. Essentially, there is no sense of "marching 
forward in time" since we can impose the boundary conditions 
at various time points. 

As a prototypical BVP we shall consider Poisson's equation in 
1D: 
\begin{equation}\label{eq:l15:poisson}
    \frac{d^2 u}{dt^2} = f(t), \quad t \in (0, 1),
\end{equation}
subject to the homogeneous Dirichlet boundary conditions 
\begin{equation*}
    u(0) = c_0, \qquad u(1) = c_1. 
\end{equation*}
This equation models many phenomenon in physics, such as 
the steady state distribution of heat in a long 
beam. Proceeding as we did for ODEs before we can rewrite 
this problem as a system of ODEs:
\begin{equation}
    \frac{d u}{d t} = z, \qquad \frac{d z}{d t} = 
    f(t),
\end{equation}
but we quickly realize that we have no initial conditions for 
$z$ and it is unclear how to impose $u(1) = u_1$. Essentially, 
our previous numerical methods are no longer applicable. 

Regardless, the mathematical ideas behind IVP solvers, 
i.e., approximating derivatives towards obtaining algebraic 
equations that we could then solve easily, remain useful. 
To be more precise, we can approximate 
\begin{equation*}
    \frac{d^2 u}{dt^2}(t) \approx 
    \frac{u(t - \Delta t) - 2 u(t) + u(t + \Delta t)}{\Delta t ^2}
    = \frac{u_{n-1} - 2 u_n + u_{n+1}}{\Delta t^2},
\end{equation*}
and substitute in \eqref{eq:l15:poisson} to obtain, with 
the boundary conditions, the system of equations 
\begin{equation}\label{eq:l15:Poisson-Fd}
   \begin{aligned}
        u_0 & = c_0 \\
        \frac{u_{n-1} - 2 u_n + u_{n+1}}{\Delta t^2} &= 
        f_n, \quad n=1, \dots, N-1, \\ 
        u_{N} &= c_1
   \end{aligned}
\end{equation}
where we used our notation from previous lectures with, 
$t_n = n \Delta t$,  
$f_n = f(t_n)$, and $N$ denoting total number of 
time steps. Note that this is a linear system of algebraic equations,
indeed we have $N+1$ equations and $N+1$ unknowns which implies 
that the system is square and may be uniquely solvable. 

Indeed, write 
\begin{equation*}
    A = \begin{bmatrix}
    1 & 0 & 0 & 0 & \dots & 0 \\ 
    1 & -2 & 1 & 0 &\dots & 0 \\
    0 & 1 & -2 & 1 & \dots & 0 \\ 
    0 & 0 & \ddots & \ddots & \ddots & \\     
    0 & \dots  & 0 &1  &-2 & 1 \\
    0 & \dots &  0 & 0 & 0 & 1
    \end{bmatrix} \in \R^{N+1 \times N+1}, 
    \qquad 
    \fb = \begin{bmatrix}
        c_0 \\ 
        f_1 \\
        f_2 \\ 
        \vdots \\
        f_{N-1} \\ 
        c_1
    \end{bmatrix} \in \R^{N+1},
\end{equation*}
as well as $\ub = (u_0, \dots, u_N) \in \R^{N+1}$. Then we 
simply need to solve 
\begin{equation*}
    A \ub = \fb.
\end{equation*}
A direct calculation reveals that the matrix $A$ is indeed 
invertible and so $\ub$ can be computed uniquely using 
standard techniques from linear algebra, e.g., 
LU or QR factorizations, iterative methods, Gaussian elimination, etc.

Notice that the fundamental difference between IVPs and 
BVPs is readily revealed. While in IVPs we compute $u_{n+1}$ 
in terms of $u_n$, or more broadly $u_0, \dots, u_{n-1}$. 
Here we compute the entire solution vector $\ub$ all at once
and so there is no sense of marching forward in time. 

Through this part of the course we primarily focus on 
second order BVPs of the form 
\begin{equation*}
    \frac{d^2 u}{dt^2} = f\left(t, u, \frac{d u}{dt} \right),
    \quad t \in (a, b)
\end{equation*}
equipped with four types of boundary conditions: 
\begin{itemize}
    \item {\it Dirichlet's B.C.}, named after Peter 
    Dirichlet:
    \begin{equation*}
        u(a) = c_a, \qquad u(b) = c_b.
    \end{equation*} 
    \item {\it Neumann's B.C.}, named after Carl Neumann: 
    \begin{equation*}
        u'(a) = c_a, \qquad u'(b) = c_b.
    \end{equation*}
    \item {\it Robin's B.C.}, named after Victor Robin:
    \begin{equation*}
        u'(a) = \alpha_a u(a) + \beta_a, \qquad 
        u'(b) = \alpha_b u(b) + \beta_b.
    \end{equation*}
    \item {\it Periodic B.C.}:
    \begin{equation*}
    u(a) = u(b), \qquad u'(a) = u'(b).     
    \end{equation*}
\end{itemize}
We will allow for the Dirichlet, Neumann, and Robin 
conditions to be mixed and matched, e.g., we have Dirichlet 
at $a$ but Robin at $b$ and so forth.

\begin{quote}
    Our goal is to construct algorithms that compute consistent 
    approximations to the true solution $u(t)$, i.e., 
    $u_n \to u(t_n)$ as $\Delta t \downarrow 0$. We shall 
    further ask for rates of convergence just as we did 
    before for IVPs.
\end{quote}

% \bibliographystyle{plainnat}
% \bibliography{references}